\section{怎样辨析单纯词和合成词}
我们在学习汉字时，可以发现汉字有这样三种类型：

一是有一定的字形和字音，却没有一定的字义，也不能独立运用。如“匍、匍、玻、璃、圆、囷”等。它们仅仅是记录语言的符号和读音的单位，这就是字。

二是有一定的字形、字音和字义，却不能独立运用。如“祖、言、肃、达、灿、维”等。它们不仅具有字的特点，而且具有最小的意义，只是不能独立运用罢了，这就叫词素，也就是具有最小意义的构词单位。

三是既有一定的字形、字音和字义，又能独立运用。如“人、水、说、笑、红、我”等。它们不仅具有字和词素的特点，而且又是最小的能独立运用的造句单位，这就是词了。

词按音节来分，有一个音节构成的单音词，有两个以上音节构成的复音词（一个汉字基本代表一个音节）；按构成一个词的词素数目来分，则有单纯词和合成词。

单纯词是指由一个词素构成的词；合成词是指由两个以上词素构成的词。在弄清字、词素和词的不同特点的前提下，我们就可以着手辨析单纯词和合成词。具体说来应注意以下几点：

1. 辨析一个词究竟由几个词素所构成。如果只由一个词素所构成，当然是单纯词，否则就是合成词。

根据这个标准，我们自然能推断出：凡是由一个词素构成的单音词，全部是单纯词，因为用一个单音节词素构成的词，字和词是一致的。例如：

山、水、田、马、牛、羊、天、地、人；我、你、他；\\
\indent
看、走、来、去；好、坏、红、蓝；一、十、个、辆；\\
\indent
的、地、了、过；和、与；从、比；才、又；啊、唉。

问题是，现代汉语以双音节词为主。我们要辨析一个双音词究竟是单纯词还是合成词，关键在于辨析它是由一个词素所构成，还是由两个词素所构成。例如：
“蝙蝠”这个词中，“蝠”字没有意义，“蝠”字也没有意义，两个字合起来才有意义，表示一种能在夜空中飞翔的哺乳动物。这就说明构成这个双音节词的每一个字，都不是词素；两个字合起来才是一个词素。这个由一个词素构成的词当然是单纯词。
“黄莺”这个词中，“黄”字有意义，“莺”字也有意义，说明它们每个字都是词素；这两个词素合起来成为一个词，表示一种身体黄色，叫声好听的鸟，而这个词的意义包含着构成它的那两个词素的意义，只是已成为一个整体了。“黄”还可以和“蜂”、“雀”、“牛”、“羊”、“金”、“豆”等构成新词；“莺”还可以同“夜”构成“夜莺”。象“黄莺”这类由两个词素构成的词，当然都是合成词了。

2. 弄清单纯词和合成词的基本类型，是辨析它们的重要手段。
汉语既以双音节词为主，而双音节的单纯词的基本类型
\newpage
\noindent
就是联绵词，所以我们只要弄清了联绵词的规律，也就掌握了双音单纯词的基本类型了。

联绵词都是双音节的，包括三种形式：

（1）双声词，即两个字的声母相同的联绵词，如：
蜘蛛（声母zh 同）、伶俐（声母l 同）、澎湃（声母p 同）、参差（声母c 同）、慷慨（声母k 同）、仿佛（声母f 同）、弥漫（声母m 同）、玲珑（声母l 同）、恍忽（声母h 同）等等，都是单纯词。

（2）叠韵词，即两个字的韵母相同的联绵词，如: 
逍遥（韵母iao同）、从容（韵母ong同）、灿烂（韵母an同）、螳螂（韵母ang同）、罗唆（韵母uo同）、薛荔（韵母i 同）、妖娆（韵母ao同）、哆嗦（韵母uo同）、依稀（韵母l 同）、翩跹（韵母lan同）、料峭（韵母lao同），等等，都是单纯词。

（3）非双声非叠韵的联绵词。这一类由两个音节合起来表示一个意义的联绵词，同双声词、叠韵词一样只包含一个词素，所以不能拆开来讲解，仍然属于单纯词。如：

玛瑙、猬孙、玻璃、萝卜、蚱蜢、蝴蝶、蝙蝠、容与、朦胧、模糊、峥嵘、圆图、刺猬、憧憬

至于合成词的基本类型就是复合式合成词。复合式合成词，一般是由两个不同的有实在意义的词素按一定的意义关系结合起来所组成。这类合成词的构成方式主要有以下五种：

（1）联合式：由意义相同、相近或者由意义相反、相对的两个词素并列起来所构成。如：

朋友、群众、领袖、路线、帮助、爱护、学习、阻塞、
\newpage
\noindent
智慧、巨大、清晰、严肃；

\noindent
呼吸，是非、矛盾、忘记、出纳、动静、始终

（2）偏正式：前一个词素修饰或限制后一个词素，一偏一正，偏在前，正在后。如：

钢笔、红旗、火车、雪白、火热、微笑、粉碎

（3）述补式：后一个词素补充、说明前一个词素，一正一偏，正在前，偏在后。如：

肃清、抓紧、立正、改善、提高、戳穿、推广

（4）动宾式：前一个词素表示动作或行为，后一个词素表示动作或行为所支配的对象。如：

碰壁、动员、司机、挂钩、裹腿、提纲、革命

（5）主谓式：前一个词素表示事物，后一个词素表示动作或性状，说明前一个词素。如：

冬至、雪崩、月蚀、国庆、民主、自觉、心疼

另外还有一种比较特殊的构词方式，就是前一个词素表示事物，后一个词素表示事物的单位，合起来组成一个词，表示一种类名。如：

诗篇、纸张、枪支、船只、文件、房间、布匹

我们明确了这几种复合式合成词的构词方式，就能够基本辨析一般的合成词了。

3. 弄清单纯词和合成词的其他类型。

单纯词除了由一个词素构成的单音词和由一个词素构成的联绵词之外，还包括译音词。译音词主要是指根据外国语言中语词的读音用汉语中跟它发音相同或近似的语音表示出来的词，也包括根据少数民族语言中的语词的读音表示出来的词。如：
\newpage
雷达、逻辑、英吉利、梵哑铃、奥林匹克、澳大利亚、布尔什维克、英特纳雄耐尔（外来语译音词）；

哈达、镭镭、葡萄、吐鲁番、乌兰浩特（少数民族语译音词）

我们了解到这一点，就可以全部辨析汉语单纯词了。
至于合成词，除了复合式构词方式之外，还有以下几种构词方式：

（1）附加式：在一个有实在意义的词素前边或后边加上一个作辅助成分的词素，合起来组成一个词。如：
①前边加辅助成分：第一、初十、小王、老师、阿姨、大李

②后边加辅助成分：法子、活儿、甜头、读者、作家、党性、海员、绿化

（2）简缩式：

①把全称词组简化成一个合成词。如：

人民代表大会——人大

中华人民共和国——中国

②选取全称中具有共性的关键性词语，用数字加以概括。如：

身体好、学习好、工作好——三好

工业、农业、科技、国防现代化——四化

我们掌握了以上所说的基本知识，并进行一些辨析单纯词和合成词的练习，就一定能准确地区分单纯词和合成词。
\newpage